% Ad Atticum 13.51 --- headnote
% ad-atticum-13-51, 24 August 45 BC, Tusculum

\textit{Cicero to Atticus, written at the Tusculan villa on the
ninth day before the Kalends of September 45 BC --- Perseus
dateline \textit{Scr. in Tusculano ix K. Sept. a. 709 (45)}. Two
sections. The first follows up on the letter Cicero composed to
Caesar about the \textit{Anticatones} (the previous letter, Att.
13.50): Atticus had asked for a copy and Cicero had forgotten to
send one; he wants Atticus to know he was not embarrassed before
him at the prospect of playing some kind of ridiculous part ---
the noun is corrupt in the manuscripts (preserved here as
\dag micillus \dag) and probably names a stock comic figure, a
\textit{little Cyrus} or \textit{little man}. He wrote, he
insists, as one peer to another \textit{[Greek: pros ison
homoion]} and without flattery \textit{[Greek: akolakeutos]} ---
the two Greek tags carrying the methodological claim about the
letter's tone.}

\textit{The second section is the usual mix: relief that Atticus's
daughter Attica is past some illness (``now at last I have
certainty''); the standing demand for full news of Tigellius's
take from Caesar; and the imminent arrival of Cicero's brother
Quintus, who is somewhere on the road and may turn up either at
Tusculum or at Atticus's house in Rome. ``I shall have to come up
to Rome myself, or he will be flying over before I get there''
--- the participial \textit{advolet} catches the harassed pace of
the summer.}
